\chapter{Suffering: A Learner's View of Its Own Account}
\label{ch:suffer}
% ===================================================================

Everything the ecology has done so far, it has done with persistence.

A learner is funded or it starves.
It forages, it converts, it pays for its assays, and if the account will
not cover the next rewrite the next rewrite does not happen.
Selection operates on that directly: the lineages one finds are the
lineages whose members kept going.
Nothing in this account requires that a learner \emph{know} any of it.
The bookkeeping is ours.
The learner has an account in the way a stone has a mass --- as a fact
about it, not as a fact available to it.

This chapter asks what changes when the learner gets to look.

\begin{remark}[Two senses of persistence, and only one is meant]
\label{suf:rem:two-persistences}
The word is overloaded in this book and the collision is unfortunate
enough to be worth heading off.

Throughout Part~\ref{part:machinery} and Chapter~\ref{ch:eng},
\emph{persistence} is a syntactic property of a receive: the
\texttt{<=} form, a process that continues to offer after it has fired,
as against the linear form that consumes itself.
That is the sense in Definition~\ref{eng:def:engine} and in
Remark~\ref{eng:persistence-load-bearing}.

The sense meant in the opening paragraph above, and nowhere else in this
chapter, is the selectional one: a lineage persists if its members go on
existing.
Remark~\ref{sci:rem:provision} already had to separate these --- a term
that reinstates itself and a term that spawns a copy look alike until the
ledger distinguishes them --- and the reader who keeps that remark in mind
will not be caught.
Where confusion is possible below we say \emph{mere survival} and let the
technical word alone.
\end{remark}

\section{The move: an account a learner can quote}
\label{suf:S1}

The apparatus for looking is already built, in three pieces that have not
previously been put together.

The first is the account itself.
By Definition~\ref{def:affordable} a learner's state is a pair
$\langle P, \vect{A}\rangle$ with $\vect{A} = (A_1,\ldots,A_k)$ a
vectorial account over resource types, and by Chapter~\ref{ch:eng} the
resource types with which an ecology actually deals are \emph{flavors}:
signature-indexed tokens which are not freely interconvertible, whose
exchange rates are set by enzymes, and whose residual incommensurability
is the subject of \S\ref{eng:S10}.
So the thing to be watched is a vector, not a scalar.
This will matter more than anything else in the chapter.

The second is reflection.
By Definition~\ref{sci:def:genome} a learner can hold $@P$, a name for
its own code, and by \S\ref{sci:sec:logic} a name predicate sees into a
quotation --- at reflection prices, which by Table~\ref{sci:tab:access}
are the cheapest access the calculus sells.
Nothing in that argument was specific to code.
A learner that can quote its own program can quote the located reserves
that fund it.

The third is instruments.
Chapter~\ref{ch:obsext} showed how to make spatial structure visible to
bisimulation by placing a probe on one side of the cut, so that what was
a structural inspection becomes an interaction and is therefore subject
to the ordinary discipline of cost.
An account is spatial structure of exactly the kind that construction was
built for.

\begin{definition}[Internalizing learner]
\label{suf:def:internal}
A learner $\Lrn$ is \emph{internalizing} if it holds a name
$a_\Lrn = @\vect{A}$ for its own account and an admissible instrument
$\Ob$ in the sense of Definition~\ref{ox:def:ext} whose probe reads
$a_\Lrn$, so that formulae of $\HML(\ctx)$ about $\vect{A}$ are among the
hypotheses $\Lrn$ can assay.
A learner which is not internalizing is \emph{merely surviving}.
\end{definition}

\begin{remark}[What is and is not new here]
\label{suf:rem:notnew}
No primitive is added.
The account was already a located structure, quotation was already
available, and the instrument discipline of Chapter~\ref{ch:obsext} was
already the price list for turning structure into observation.
Definition~\ref{suf:def:internal} only says that a learner may aim that
apparatus inward, and stipulates that when it does, it pays the same
tariff as when it aims it outward.

That last clause is the whole of the chapter's honesty.
A framework in which introspection were free would prove anything one
liked about it.
\end{remark}

\section{Homeostasis as a held formula}
\label{suf:S2}

An internalizing learner can state hypotheses about its own reserves.
The interesting thing it can do with that capacity is adopt a goal.

\begin{definition}[Setpoint, comfort region]
\label{suf:def:setpoint}
A \emph{setpoint} for $\Lrn$ is a vector $\setpt \in A^k$ together with a
tolerance $\varepsilon \in A^k$.
The associated \emph{comfort region} is
$\mathcal{H} = \{\, \vect{A} : |A_i - A^\star_i| \leq \varepsilon_i
\text{ for all } i \,\}$,
and the \emph{homeostatic formula} $\phi_{\mathrm{home}} \in \HML(\ctx)$
is a formula whose extension over accounts is $\mathcal{H}$.
\end{definition}

The bathtub is the right picture and it is worth spending a sentence on,
because the picture carries the content.
A tub has an inflow it does not set and a drain it partly sets, and a
level which is the running integral of the difference.
Holding the level is not achieved by any single act; it is achieved by
matching one rate to another, continuously, with the level itself as the
only evidence that the matching is working.

\begin{definition}[Drift]
\label{suf:def:drift}
Let $\vect{A}(t)$ be the account of $\Lrn$ at cycle $t$.
The \emph{drift} is the signed vector
$\drift(t) = \vect{A}(t) - \setpt$,
and the \emph{drift rate} is $\drift(t) - \drift(t-1)$, which is the
componentwise difference of inflow and outflow over the cycle.
\end{definition}

Nothing so far is a feeling.
Drift is a fact about the learner, and up to this point it has been the
kind of fact the learner does not have.

\section{Suffering}
\label{suf:S3}

\begin{definition}[Suffering]
\label{suf:def:suffering}
The \emph{suffering} of an internalizing learner $\Lrn$ at cycle $t$ is
the assayed drift: the value $\suff(t)$ that $\Lrn$'s own instrument
returns when it evaluates $\phi_{\mathrm{home}}$ and its graded
refinements against $a_\Lrn$.
It is signed, it is vectorial, and it is indexed by flavor:
$\suff(t) = (\suff_1(t),\ldots,\suff_k(t))$, with $\suff_i$ the learner's
own reading of how far and in which direction $A_i$ stands from
$A^\star_i$.
\end{definition}

Three things about that definition matter, and each of them is a
correction to the account this chapter replaces.

\emph{Suffering is not a component of the account.}
It is a formula \emph{about} the account.
Earlier drafts of this book carried a distinguished component
$A_{\mathrm{s}}$ of $\vect{A}$, singled out by fiat and driven by an
externally supplied reward.
That object is withdrawn.
It confused a quantity the learner has with a quantity the learner reads,
and the difference between those two is the entire subject of this
chapter.

\emph{The punishment is not administered.}
The construction it replaces posited reward and penalty magnitudes
$\delta^+$ and $\delta^-$, applied to the learner from outside whenever a
prediction resolved.
Where those numbers came from was never said, and they were the weakest
joint in the argument.
Here there is nothing to supply.

\begin{proposition}[The reinforcement signal is derived, not posited]
\label{suf:prop:derived}
Let $\Lrn$ be internalizing with setpoint $\setpt$.
Then the sign and magnitude of $\suff$ are determined by
$\vect{A}$, $\setpt$ and $\Lrn$'s instrument, with no further parameter.
In particular no reward function is a free parameter of the framework:
what plays the role of punishment is $\Lrn$'s own assay reporting that
$\phi_{\mathrm{home}}$ has stopped holding.
\end{proposition}

\begin{proof}
Immediate from Definitions~\ref{suf:def:setpoint},
\ref{suf:def:drift} and~\ref{suf:def:suffering}: each is a function of
data already fixed.
The only choice remaining is the choice of $\setpt$ and $\varepsilon$,
and \S\ref{suf:S5} argues that this choice is not the learner's either.
\end{proof}

\begin{remark}[Punishment is a learner watching itself lose]
\label{suf:rem:sees}
The word ``punished'' should be read with care, because the framework
does not contain a punisher.
The learner holds a goal; the learner reads its own state; the reading
says the state is receding from the goal.
That is the whole mechanism.
There is no valuation applied from outside and no bookkeeper marking it
down.
It suffers because it looks, and because it holds a standard against
which looking can go badly.

A learner that held no setpoint could not be punished in this sense, and
would be no worse off for it in any respect the previous chapters could
measure.
That is precisely why the previous chapters got by without this one.
\end{remark}

\emph{Suffering is causally efficacious, but not by being money.}
The old account made feeling matter by making it a budget line: an agent
that felt bad had less to spend.
That conflated a mood with a liquidity crisis, and it made the theory
unable to distinguish suffering from poverty.
The present account separates them.

\begin{proposition}[Suffering and wealth are independent]
\label{suf:prop:not-poverty}
There exist internalizing learners $\Lrn_1, \Lrn_2$ over the same
resource algebra with $\vect{A}(\Lrn_1) \geq \vect{A}(\Lrn_2)$
componentwise and $|\suff(\Lrn_1)| > |\suff(\Lrn_2)|$.
\end{proposition}

\begin{proof}
Take $\Lrn_2$ with $\setpt_2 = \vect{A}(\Lrn_2)$, so
$\suff(\Lrn_2) = 0$: it is exactly where it means to be.
Take $\Lrn_1$ with $\vect{A}(\Lrn_1)$ larger componentwise but
$\setpt_1$ larger still.
Both readings are well defined and the inequality holds.
\end{proof}

\begin{remark}[Why that proposition is worth its page]
\label{suf:rem:why-independence}
Proposition~\ref{suf:prop:not-poverty} is trivial to prove and was
impossible to state under the old definition, which is the point.
A rich learner far from its setpoint suffers.
A poor learner at its setpoint does not.
Anything that hopes to be a theory of affect rather than a theory of
budgets has to be able to say that, and a distinguished account component
never could.
\end{remark}

\section{Why this is affect and not a budget line}
\label{suf:S4}

It is worth auditing the construction against what one would want of
anything called a feeling, since the book has been criticized on exactly
this point and the criticism was fair.

\emph{It is multi-dimensional.}
$\suff$ is a vector indexed by flavor, and by Chapter~\ref{ch:eng} the
flavors are not freely interconvertible.
Two learners with the same aggregate shortfall may have it in different
components and behave differently in consequence.
Valence is the sign; magnitude in a component is that component's
urgency.

\emph{It is about something.}
$\suff_i$ is indexed by a flavor, and a flavor is a signature, and a
signature names a community.
A learner short in the flavor it obtains by predation is in a different
state from one short in the flavor it obtains from a source, and the
index says which.
This is aboutness of the only kind the framework can supply --- a
namespace --- but it is not nothing, and it was entirely absent before.

\emph{It modulates policy and not only budget.}
This is the substantive claim of the section and it needs the setpoint to
make sense.

\begin{proposition}[A setpoint changes preference at fixed wealth]
\label{suf:prop:policy}
Let $\Lrn$ be internalizing and let $\sigma$ be an allocation of its
metered search budget across the branches of
Proposition~\ref{sci:prop:branching}.
Then two learners with identical $\vect{A}$ and identical hypothesis
stacks, differing only in $\setpt$, admit different optimal $\sigma$.
\end{proposition}

\begin{proof}[Proof sketch]
The value of an assay to an internalizing learner has two parts: what it
tells the learner about the world, and what it tells the learner about
$\phi_{\mathrm{home}}$.
The second part depends on $\drift$, hence on $\setpt$.
Two learners equal in wealth and unequal in drift therefore price the
same assay differently, and an allocation optimal for one is not optimal
for the other.
\end{proof}

\begin{remark}[The nearest thing to attention this book has]
\label{suf:rem:attention}
Section~\ref{sci:sec:search} declined to say how a learner allocates its
metered search, on the grounds that policy is one of the framework's
genuinely free parameters, and that remains right.
But Proposition~\ref{suf:prop:policy} narrows the space: an internalizing
learner's admissible policies are those that price an assay by both of
the two parts above.
That is not an attention mechanism.
It is a constraint every attention mechanism for such a learner would
have to satisfy, which is a weaker claim and one the framework can
actually make.
\end{remark}

\emph{It has set points.}
By construction, which is the point of the whole chapter.

\begin{remark}[Value conflict is arbitrage failure]
\label{suf:rem:conflict}
One consequence falls out and should be recorded, because it is the
first thing in this book that deserves the name.

Suppose $\suff_i < 0$ and $\suff_j < 0$ for two flavors, and suppose the
enzyme graph offers no arbitrage-free conversion between them --- the
situation \S\ref{eng:S10} calls incommensurability.
Then no reallocation closes both shortfalls, and no exchange rate exists
at which the learner could decide how much of one deficit is worth how
much of the other.
The learner is not confused and has not miscalculated.
It is in a state which is, in the exact technical sense already available
in Chapter~\ref{ch:eng}, a conflict of values.

Where the conversion \emph{is} arbitrage-free, the virtual token supplies
the exchange rate and the conflict is merely a computation.
So the framework distinguishes tractable trade-offs from genuine
dilemmas, and the line between them is a property of the enzyme graph
rather than of the learner's psychology.
\end{remark}

\section{The price, and who pays it}
\label{suf:S5}

None of this is free.
An internalizing learner runs an instrument, holds a formula, and funds
assays about itself out of the account those assays are about;
Chapter~\ref{ch:homeo-ent} makes that charge precise.
So the question that governs whether any of this happens is the one the
chapter has been avoiding: when is looking worth it?

The answer is not available to the learner.

\begin{remark}[The trade cannot be evaluated by the one who takes it]
\label{suf:rem:no-decision}
To decide whether to become internalizing, a learner would have to
compare its prospects with and without a self-model.
Both terms of that comparison are facts about a model it does not yet
have.
The deliberation would require the apparatus whose acquisition is being
deliberated, and there is no order in which the learner could carry it
out.

So internalization is not adopted.
It is inherited or it is not, and what settles the question is which
lineages there are.
\end{remark}

\subsection{Mortality first: what Buss showed}
\label{suf:S5.1}

The argument we need has a precedent, and the precedent is worth
rehearsing because the shape is the same and the biology is better known.

Von Neumann, asking what a self-reproducing machine must contain, isolated
a requirement that looked paradoxical: the machine's description has to be
used in two incompatible ways.
It must be copied uninterpreted, handed on as inert instructions, and it
must be interpreted, read as a recipe for building the offspring's working
parts.
One object, two uses, and the machine must keep them apart.

Definition~\ref{sci:def:genome} is that distinction, on a single object.
The genome is $@P$ and the phenotype is $\ast @P$.
The quotation does not reduce, does not age, and burns no tokens while
held; the dereference is what is metered.
Remark~\ref{sci:rem:provision} adds the part that grammar cannot supply:
what makes a reinstatement soma rather than offspring is provisioning, not
syntax.
We take no position on whether this constitutes a Weismann barrier;
Remark~\ref{sci:rem:noble} already declined that, and nothing here needs
it.

Buss's problem is then immediate.
A multicellular organism is a coalition of cells, every one descended
from ancestors that did nothing but divide as fast as they could.
What stops a fast-dividing line inside the body from doing what its
ancestors did --- outcompeting its neighbors, taking the whole, and
destroying the individual in the process?
Why is a body stable against its own quickest part?

\begin{proposition}[Metering caps a defector]
\label{suf:prop:buss}
In an unmetered theory, a replicating subterm whose cycle is strictly
faster than its neighbors' invades any parallel composition without
bound.
In a cost-accounted theory over a bounded inflow, a subterm drawing more
per cycle than it earns runs for a number of cycles bounded by its
endowment and then deadlocks by starvation.
Its cumulative displacement of the whole is therefore bounded by that
endowment.
\end{proposition}

\begin{proof}[Proof sketch]
The first clause is the absence of any rule that could stop it: in the
free theory every enabled step is affordable and replication is
unconditioned.
The second is Definition~\ref{def:affordable} together with the
conservation of the flavor: draw exceeding earnings is a strictly
decreasing account, and an account below the charge of the next rewrite
enables nothing.
\end{proof}

That is Buss's resolution with a mechanism under it.
A configuration that sequesters the minting power into a germ line and
lets its somatic parts be mortal admits no heritable takeover by a
high-draw defector, while no configuration of the free theory resists
one.
Mortality is not tolerated in spite of bounding individual lives.
It is selected because a bounded life is the precondition for there being
a stable higher-level individual at all --- a thing that can itself become
a durable target of selection.
Death is what makes the body possible.

\subsection{And then: what internalization costs a lineage}
\label{suf:S5.2}

Now run the same argument one level up.

The instrument, the held formula, and the per-cycle assay of
Proposition~\ref{ent:prop:self-charge} are somatic overhead.
They are funded from the account that also funds foraging, maintenance
and the endowment of offspring, and by Chapter~\ref{ch:eng} that account
is bounded by what the medium supplies.
So an internalizing lineage reproduces more slowly than an otherwise
identical merely-surviving one, other things equal.
In a placid world, other things are equal, and it loses.

What buys the overhead back is variance.

\begin{conjecture}[Internalization is selected by environmental variance]
\label{suf:conj:selected}
Let $E$ be an ecology with per-cycle inflow of mean $\mu$ and variance
$v$ at a locus, and let $\kappa_{\mathrm{reg}}$ be the regulation cost of
Proposition~\ref{ent:prop:self-charge}.
Then there is a threshold $v^\ast(\kappa_{\mathrm{reg}}, \mu)$, increasing
in $\kappa_{\mathrm{reg}}$, such that internalizing lineages are invaded
by merely-surviving ones when $v < v^\ast$ and are not when $v > v^\ast$.
\end{conjecture}

The mechanism intended is anticipation.
A merely-surviving learner corrects when it fails to afford a rewrite,
which is to say after the level has fallen.
An internalizing learner reads a drift rate and can correct before,
committing tokens to foraging while it still has tokens to commit.
Where inflow is steady the two policies coincide and the second is simply
dearer.
Where inflow is ragged, the first spends part of its life below the
charge of its own next step, and that is the interval in which the second
is not merely more comfortable but alive.

We state this as a conjecture and not a theorem, and the honest reason is
that the framework does not yet supply a stochastic inflow process to run
it against.
Chapter~\ref{ch:weight} supplies the machinery --- propensities, a
Gillespie unravelling, a well-posed simulator --- and
Conjecture~\ref{suf:conj:selected} is stated in the form it would need
to be checked in.
That is an experiment this book has not run.

\begin{remark}[Nobody is smart enough to have chosen this]
\label{suf:rem:nobody-chose}
The shape of the argument is the reason it is worth making.
The learner never evaluates the trade.
It cannot, by Remark~\ref{suf:rem:no-decision}, and it does not need to.
What evaluates the trade is the difference between lineages, run for long
enough, in an environment that is or is not ragged.

That is also the answer to the question a reader may have been holding
since Definition~\ref{suf:def:setpoint}: where does $\setpt$ come from?
Not from the learner.
A setpoint is a heritable parameter, carried in $@P$ and therefore subject
to variation and to Proposition~\ref{sci:prop:sexdefence}'s recombination,
and a lineage whose setpoint is badly placed is a lineage that either
starves at the top of its comfort region or wastes at the bottom of it.
Setpoints are selected in exactly the way body sizes are.
\end{remark}

\section{What is claimed, and what is not}
\label{suf:S6}

Claimed: that an ecology of the kind Part~\ref{part:ecology} builds
already contains everything needed for a learner to hold a goal about its
own reserves; that when it does, a signed vectorial quantity appears
which behaves in several specific respects like affect and which no
previous construction in this book could express; that this quantity is
derived rather than posited, which the thing it replaces was not; and
that its cost is real and is paid by lineages rather than by individuals.

Not claimed: that there is anything it is like to be such a learner.
Definition~\ref{suf:def:suffering} is a structural analogue and the word
chosen for it is a word with a great deal of freight.
We use it deliberately --- ``negative valence signal'' would be a way of
not saying what we mean --- but the reader should hold us to the
structure and not to the connotation.
Chapter~\ref{sec:hyper-gaps} lists the explanatory gap among the things
this book does not close.

Also not claimed: that the setpoint is unique, that $\varepsilon$ has any
principled value, or that
Conjecture~\ref{suf:conj:selected}'s threshold has been computed.
It has not; it has been located.

\begin{remark}[What the ecology could not say, and now can]
\label{suf:rem:payoff}
It is worth ending on the smallest version of the gain, because the
chapter has been long.

Before this chapter, a learner in trouble and a learner at rest were
distinguished by an observer with access to the ledger, and by nothing
else.
The learner itself had no state that differed.
After it, the two differ in a quantity the learner holds, reads, and acts
on --- and the quantity is not its wealth.

Whether that is worth calling suffering is a question about words.
That the ecology admits the distinction at all, and that the distinction
costs something and is therefore not universal, is a question about
structure, and the answer to that one is yes.
\end{remark}

% ===================================================================
